\section{Methodology}
\label{sec:methodology}

We present \textbf{FIG-MAC}, a collaborative multi-agent framework designed to automate scientific hypothesis generation by simulating the reasoning process of human research teams. As shown in Figure~\ref{fig:framework}, FIG-MAC architecture comprises three synergistic modules: (i) a \emph{fine-grained inspiration graph} (FIG) that provides a structured knowledge substrate, (ii) a \emph{hybrid retrieval} mechanism that bridges semantic search with structural reasoning, and (iii) a \emph{multi-agent workflow} that drives iterative hypothesis refinement through adversarial peer review.

\subsection{Knowledge Base: Fine-grained Inspiration Graph}
\label{sec:graph_construction}

The effectiveness of hypothesis generation hinges on the granularity of knowledge. Unlike traditional citation graphs that treat papers as atomic nodes, we model the scientific literature space as a heterogeneous graph $\mathcal{G} = (\mathcal{V}, \mathcal{E})$ that captures the internal logical skeleton of research.

\paragraph{Node and Relation Modeling.}
We define three primary entity types in $\mathcal{V}$: \textbf{Papers ($P$)}, \textbf{Research Questions ($Q$)}, and \textbf{Solutions ($S$)}. Each Paper node is initialized with semantic embeddings, while $Q$ and $S$ nodes represent structured elements extracted from the text. Relationships in $\mathcal{E}$ are categorized into hierarchical (e.g., \texttt{has\_rq}) and cross-document (e.g., \texttt{inspired}). The \texttt{inspired} relation acts as the engine for innovation, representing pathways where one study's solution potentially addresses another's open question.

\paragraph{Construction Pipeline.}
The FIG is constructed via a three-stage automated pipeline. First, an \textbf{Information Extraction} module identifies core semantic units using pattern matching and LLM-based parsing. Second, \textbf{Entity Embedding} maps these units into a high-dimensional vector space using a pre-trained encoder. Finally, \textbf{Relationship Establishment} identifies cross-paper links by combining vector similarity filtering (threshold $\tau = 0.5788$) with LLM-based classification (Qwen3-8B, 0.95 accuracy) and heuristic enrichment. The resulting graph stores 18.8k unique entities and 1.8M edges in a dual storage architecture, enabling simultaneous structural and semantic queries.

\subsection{Hybrid Retrieval and Evidence Synthesis}
\label{sec:hybrid_retrieval}

To provide agents with comprehensive context, FIG-MAC employs a hybrid retrieval mechanism $\mathcal{R}(q)$ that complements semantic "flesh" with structural "skeleton."

\paragraph{Dual Retrieval Streams.}
Given a query $q$, the system executes \textbf{Vector Retrieval} $\mathcal{R}_V(q)$ across multiple attributes (abstracts, methods) to gather diverse textual evidence. Simultaneously, \textbf{Graph Retrieval} $\mathcal{R}_G(q)$ maps the query to seed nodes in $\mathcal{G}$ and performs multi-hop discovery to find inspiration paths $\pi = (q_i, s_j, p_k)$. These paths are ranked based on a composite score:
\begin{equation}
    \text{score}(\pi) = \mathbf{w}_{\pi}^\top \phi(\pi, q)
\end{equation}
where $\phi(\pi, q)$ captures semantic relevance and structural confidence of path segments.

\paragraph{Skeleton-Flesh Integration.}
Our system synthesizes these streams using a "Skeleton-Flesh" strategy. Inspiration paths serve as the structural framework (the skeleton) for the hypothesis, while node-specific technical details (the flesh) are enriched via vector-retrieved evidence. This integration ensures that the generated hypotheses are not only novel but also technically grounded and traceable to specific literature.

\subsection{Inspiration Path Generation via Link Prediction}
\label{sec:link_prediction}

To uncover innovation pathways not explicitly mentioned in literature, we adopt a Relation Graph Convolutional Network (RGCN) to update node representations:
\begin{equation}
    \mathbf{h}_i^{(l+1)} = \sigma \left( \sum_{r \in \mathcal{R}} \sum_{j \in \mathcal{N}_r(i)} \frac{1}{c_{i,r}} \mathbf{W}_r^{(l)} \mathbf{h}_j^{(l)} + \mathbf{W}_0^{(l)} \mathbf{h}_i^{(l)} \right)
\end{equation}
Potential links are decoded using a DistMult scorer $s(u, r, v)$, and high-confidence latent edges are extracted via greedy beam search. This allows FIG-MAC to perceive potential research gaps that may have been overlooked by human reviewers.

\subsection{Hierarchical Context Management}
\label{sec:context}

To optimize information retention under token constraints, FIG-MAC employs a hierarchical management system rooted in information theory. 

\paragraph{Agent-Level Memory.}
Each agent maintains a dual-memory structure $M_i = (M_i^p, M_i^c)$. To maintain relevance, a subset $M_i^{c'} \subseteq M_i^c$ is selected via:
\begin{equation}
    M_i^{c'} = \arg\max_{\mathcal{M}'} \sum_{\mathbf{m} \in \mathcal{M}'} \mathbf{w}^\top \phi(\mathbf{m}, s_i, t)
\end{equation}
The feature vector $\phi$ captures recency $\exp(-\lambda \Delta t)$, semantic relevance $\cos(\text{embed}(\mathbf{m}), s_i)$, and LLM-judged quality.

\paragraph{Workflow-Level Context.}
The system maintains a dynamic context graph $\mathcal{G}_C$ where nodes represent phase results and edges denote information dependencies. Context aggregation follows PageRank centrality: $C_{\text{wf}}^{(t)} = \text{Select}(\mathcal{G}_C, L_{\max})$, ensuring that the most influential evidence is preserved across collaboration phases.

\subsection{Collaborative Execution and Iteration}
\label{sec:workflow}

The transition logic of FIG-MAC ensures that the hypothesis generation process is both parallelizable and self-correcting.

\paragraph{Parallel Analysis.}
For independent analysis tasks (e.g., technical vs. ethical review), agents execute concurrently. The final analysis report $R_{\text{Analysis}}$ is a weighted aggregation:
\begin{equation}
    R_{\text{Analysis}} = \sum_{i \in \text{Agents}} \omega_i \cdot \text{normalize}(R_i)
\end{equation}
where $\omega_i$ is determined by the agent's confidence score, achieving linear speedup relative to sequential execution.

\paragraph{State-Driven Refinement.}
Collaboration is governed by a state machine $\mathcal{SM}$ guiding the system through \textbf{LITERATURE}, \textbf{IDEATION}, \textbf{ANALYSIS}, and \textbf{REVIEW}. Crucially, FIG-MAC incorporates a conditional loop for self-improvement. After the REVIEW phase, the transition function $\mathcal{T}$ determines the next state based on quality threshold $\theta$:
\begin{equation}
    \mathcal{T}(\text{REV}, c) = 
    \begin{cases} 
        \text{POLISH} & \text{if } q_{ov} \geq \theta \lor i \geq I_{\max} \\
        \text{REVISE} & \text{otherwise}
    \end{cases}
\end{equation}
During the \textbf{REVISE} phase, the leader agent addresses identifies weaknesses systematically. This process repeats until convergence or reaching $I_{\max}$.

\paragraph{Quality Evaluation.}
The final output is assessed via an integrated metric $q_{int} = \lambda \bar{q}_{int} + (1-\lambda) \bar{q}_{ext}$, balancing internal peer review scores with external assessments of feasibility and novelty. This ensures FIG-MAC produces hypotheses that are technically sound and publication-ready.
